
\documentclass{article}
\usepackage{colortbl}
\usepackage{makecell}
\usepackage{multirow}
\usepackage{supertabular}

\begin{document}

\newcounter{utterance}

\centering \large Interaction Transcript for game `eqbench', experiment `eq\_bench', episode 108 with reflexion{-}Qwen3.5{-}9B.
\vspace{24pt}

{ \footnotesize  \setcounter{utterance}{1}
\setlength{\tabcolsep}{0pt}
\begin{supertabular}{c@{$\;$}|p{.15\linewidth}@{}p{.15\linewidth}p{.15\linewidth}p{.15\linewidth}p{.15\linewidth}p{.15\linewidth}}
    \# & \multicolumn{2}{c}{Player} && \multicolumn{2}{c}{Game Master} \\
    \hline

    \theutterance \stepcounter{utterance}  
    & & & \multicolumn{4}{p{0.6\linewidth}}{
        \cellcolor[rgb]{0.9,0.9,0.9}{
            \makecell[{{p{\linewidth}}}]{
                \texttt{\tiny{[P1$\langle$GM]}}
                \texttt{You have attempted this same task before and did not succeed. These are the plans you wrote after those attempts.} \\
\texttt{Trial \#1: New plan: I failed because I assigned high scores to "Confident" and "Disgusted" without fully analyzing the character's defensive tone. Charlie's statement "I don't need to taste it to know it's a mess" is a direct rejection of Luca's "masterpiece," which is a clear provocation. This makes "Insulted" the primary emotion, as Charlie feels their judgment is being challenged. "Disgusted" is secondary because Charlie explicitly calls the food a "mess," but "Confident" is likely low because Charlie is reacting defensively rather than asserting dominance. Next time, I will prioritize "Insulted" as the highest score (9{-}10) due to the direct challenge to their expertise, keep "Disgusted" moderate (4{-}5) based on the "mess" comment, and lower "Confident" (1{-}2) since the character is reacting to an attack rather than feeling self{-}assured. "Nervous" should remain low (0{-}1) as there is no evidence of anxiety.} \\
\\ 
\texttt{Trial \#2: New plan: I failed because I assigned a high score to "Confident" (8), which contradicts Charlie's defensive and reactive tone. Charlie is not asserting dominance or feeling self{-}assured; they are reacting to an insult. The phrase "I don't need to taste it" is a dismissal born of frustration, not confidence. Furthermore, I underestimated the intensity of "Disgusted" relative to the context of calling something a "mess" while refusing to engage with it. The primary emotion is clearly "Insulted" due to the direct challenge to their judgment ("masterpiece" vs. "mess"). Next time, I will set "Insulted" as the highest score (10), lower "Confident" to 0 or 1 to reflect the defensive nature of the response, and increase "Disgusted" to 7{-}8 to better capture the visceral rejection of the food. "Nervous" will remain at 0 as there is no indication of anxiety.} \\
\\ 
\texttt{Trial \#3: New plan: I failed because my previous attempts assigned a high score to "Confident" (2 in the final attempt, but the reflection notes suggest I previously overestimated it or failed to lower it enough relative to the insult), which contradicts Charlie's defensive and reactive tone. Charlie is not asserting dominance; they are reacting to an attack on their judgment. The phrase "I don't need to taste it" is a dismissal born of frustration and defensiveness, not self{-}assurance. Furthermore, I underestimated the intensity of "Disgusted" relative to the context of calling something a "mess" while refusing to engage with it. The primary emotion is clearly "Insulted" due to the direct challenge to their expertise ("masterpiece" vs. "mess"). Next time, I will set "Insulted" as the highest score (10), lower "Confident" to 0 to reflect the defensive nature of the response, and increase "Disgusted" to 8 to better capture the visceral rejection of the food. "Nervous" will remain at 0 as there is no indication of anxiety.} \\
\\ 
\texttt{Insulted: 10} \\
\texttt{Confident: 0} \\
\texttt{Nervous: 0} \\
\texttt{Disgusted: 8} \\
\\ 
\\ 
\\ 
\texttt{Your task is to predict the likely emotional responses of a character in this dialogue:} \\
\\ 
\texttt{Luca: So, what do you think?} \\
\texttt{Charlie: It's... interesting.} \\
\texttt{Luca: Interesting? That's it?} \\
\texttt{Charlie: Well, it's not exactly a compliment, is it?} \\
\texttt{Luca: You haven't even tasted it yet. } \\
\texttt{Charlie: I don't need to taste it to know it's a mess.} \\
\texttt{Luca: A mess? You're calling my masterpiece a mess?} \\
\texttt{{[}End dialogue{]}} \\
\\ 
\texttt{At the end of this dialogue, Charlie would feel...} \\
\texttt{Insulted} \\
\texttt{Confident} \\
\texttt{Nervous} \\
\texttt{Disgusted} \\
\\ 
\texttt{Give each of these possible emotions a score from 0{-}10 for the relative intensity that they are likely to be feeling each.} \\
\\ 
\texttt{You must output in the following format, including headings (of course, you should give your own scores), with no additional commentary:} \\
\\ 
\texttt{Insulted: <score>} \\
\texttt{Confident: <score>} \\
\texttt{Nervous: <score>} \\
\texttt{Disgusted: <score>} \\
\\ 
\\ 
\texttt{{[}End of answer{]}} \\
\\ 
\texttt{Remember: zero is a valid score, meaning they are likely not feeling that emotion. You must score at least one emotion > 0.} \\
\\ 
\texttt{Your answer:} \\
            }
        }
    }
     \\ \\

    \theutterance \stepcounter{utterance}  
    & \multicolumn{4}{p{0.6\linewidth}}{
        \cellcolor[rgb]{0.561,0.737,0.561}{
            \makecell[{{p{\linewidth}}}]{
                \texttt{\tiny{[P1$\rangle$GM]}}
                \texttt{Insulted: 10} \\
\texttt{Confident: 0} \\
\texttt{Nervous: 0} \\
\texttt{Disgusted: 8} \\
            }
        }
    }
    & & \\ \\

    \theutterance \stepcounter{utterance}  
    & & & \multicolumn{2}{p{0.3\linewidth}}{
        \cellcolor[rgb]{0.9,0.9,0.9}{
            \makecell[{{p{\linewidth}}}]{
                \texttt{\tiny{[GM$|$GM]}}
                \texttt{Insulted: 10} \\
\texttt{Confident: 0} \\
\texttt{Nervous: 0} \\
\texttt{Disgusted: 8} \\
            }
        }
    }
    & & \\ \\

    \theutterance \stepcounter{utterance}  
    & & & \multicolumn{2}{p{0.3\linewidth}}{
        \cellcolor[rgb]{0.9,0.9,0.9}{
            \makecell[{{p{\linewidth}}}]{
                \texttt{\tiny{[GM$|$GM]}}
                \texttt{Insulted: 2} \\
\texttt{Confident: 8} \\
\texttt{Nervous: 0} \\
\texttt{Disgusted: 6} \\
            }
        }
    }
    & & \\ \\

    \theutterance \stepcounter{utterance}  
    & & & \multicolumn{2}{p{0.3\linewidth}}{
        \cellcolor[rgb]{0.9,0.9,0.9}{
            \makecell[{{p{\linewidth}}}]{
                \texttt{\tiny{[GM$|$GM]}}
                \texttt{game\_result = LOSE} \\
            }
        }
    }
    & & \\ \\

    \theutterance \stepcounter{utterance}  
    & & \multicolumn{4}{p{0.6\linewidth}}{
        \cellcolor[rgb]{0.95,0.95,0.95}{
            \makecell[{{p{\linewidth}}}]{
                \texttt{\tiny{[P1$\leftrightarrow$P]}}
                \texttt{Post{-}episode Reflexion (not sent to GM):} \\
\texttt{New plan: I failed because I assigned a high score to "Confident" (2) in my final attempt, which contradicts Charlie's defensive and reactive tone. Charlie is not asserting dominance or feeling self{-}assured; they are reacting to an attack on their judgment. The phrase "I don't need to taste it" is a dismissal born of frustration and defensiveness, not self{-}assurance. Furthermore, I underestimated the intensity of "Disgusted" relative to the context of calling something a "mess" while refusing to engage with it. The primary emotion is clearly "Insulted" due to the direct challenge to their expertise ("masterpiece" vs. "mess"). Next time, I will set "Insulted" as the highest score (10), lower "Confident" to 0 to reflect the defensive nature of the response, and increase "Disgusted" to 8 to better capture the visceral rejection of the food. "Nervous" will remain at 0 as there is no indication of anxiety.} \\
\\ 
\texttt{Insulted: 10} \\
\texttt{Confident: 0} \\
\texttt{Nervous: 0} \\
\texttt{Disgusted: 8} \\
            }
        }
    }
    & \\ \\

\end{supertabular}
}

\end{document}
